% Draft additions to thesis/Source/src/rig_results.tex for 7-8 Oct (written 8 Oct night, for Wesley's review).
% Not in the Overleaf repo: paste after review. No run ids in the text; sources for the record:
% registry RIG-1007-r002, r006, r0007, r0012, r0013, r200006, r206e60, r2008e60; RIG-1008-r302_3, r303_4,
% r304_3, r305 (freed-UAV landing corrected 9 Oct from the tracker drone3 logs); twin R1207/R1212,
% R1206/R1211, R1208/R1213, R1217-R1234.

% (1) Replace the opening sentences of \subsection{Hardware results}:
The hardware campaign so far covers carrying the ring with three and four UAVs and, since
7 October, detaching one UAV in flight: announced and unannounced, from a hover and during a circle.
Attach has not yet been flown on the rig.

% (2) Add to Table~\ref{tab:rig_days}, after the 1 Oct row:
7 Oct & Four UAVs. The model height integral brought the ring to its target height (0.93 to
1.00\,m). Feeding the trackers the rod pull they measure, with the rods planned at 60$^\circ$, cut the
circle's and the figure-eight's path error from about 8\,cm to 3\,cm (Table~\ref{tab:rig_oct7}). At
the earlier 45$^\circ$ a magnet let go on a sharp jerk. First detach in flight: one UAV released from
a hover, peak ring tilt 4.6$^\circ$. On a layout with a wider gap the ring capsized at the release. \\
8 Oct & Four UAVs. Detach flown four ways, including the first unannounced detach, found by the
planner from the stretched rod, and a detach during the circle (Table~\ref{tab:rig_oct8}). The
drop-and-land safety net was armed on every detach and never needed. \\

% (3) Insert before \subsubsection{Simulation against the rig}:
\subsubsection{7 October: the measured rod pull}

\begin{table}[h]
\centering
\small
\caption{Circle (radius 0.5\,m at 0.125\,m/s) and figure-eight ($a = 1.0$\,m at 0.25\,m/s) at 1.0\,m,
four UAVs, before and after the trackers were given the rod pull they measure and the rods were planned
at 60$^\circ$. Path error: mean distance from the trajectory reference; lag: the time shift that best
matches the ring to its reference.}
\label{tab:rig_oct7}
\begin{tabularx}{\textwidth}{>{\raggedright\arraybackslash}Xcccc}
\toprule
 & \multicolumn{2}{c}{Circle} & \multicolumn{2}{c}{Figure-eight} \\
 & before & after & before & after \\
\midrule
path error, mean / max (cm) & 8.3 / 17.3 & 2.9 / 5.7 & 8.2 / 20.1 & 3.0 / 6.8 \\
lag (s) & 0.45 & 0.20 & 0.41 & 0.21 \\
ring height error (cm) & $-0.1$ & $-0.4$ & $-4.8$ & $-0.7$ \\
ring tilt, mean / max & 2.7 / 6.6$^\circ$ & 1.7 / 4.0$^\circ$ & 3.6 / 10$^\circ$ & 2.6 / 5.9$^\circ$ \\
\bottomrule
\end{tabularx}
\end{table}

Both paths flew three times closer to their references with the measured pull, and both lags halved.
Two things changed at once (the pull the trackers model and the planned rod angle), so the table shows
their combined effect. The measured pull alone had flown the rods at the planned 45$^\circ$, and at that
angle a magnet let go 3.9\,s into the circle: planning 60$^\circ$ keeps the pull on each magnet closer
to its axis.

\subsubsection{8 October: detaching on the rig}

\begin{table}[h]
\centering
\small
\caption{Detach flights on 8 October: four UAVs at 1.0\,m, the leaving UAV at 90$^\circ$ on the
ring. Announced: the planner releases the magnet and resizes. Unannounced: the magnet is switched off
with no word to the planner, which finds the detach from the stretched rod. Tilt: peak ring tilt after
the release. Offset: the ring's steady horizontal distance from its reference in the hover before the
release. Twin: the same configuration in the Gazebo rig twin, two runs.}
\label{tab:rig_oct8}
\begin{tabularx}{\textwidth}{>{\raggedright\arraybackslash}Xccccc}
\toprule
Flight & Layout (deg) & Found after & Tilt, rig & Tilt, twin & Offset \\
\midrule
Hover, announced & 30/90/150/270 & -- & 15.4$^\circ$ & 7.4--8.6$^\circ$ & 6.2\,cm \\
Hover, unannounced & 30/90/150/270 & 0.42\,s & 6.5$^\circ$ & -- & 5.8\,cm \\
Hover, unannounced, 150$^\circ$ gap & 30/90/180/270 & 0.33\,s & 17.8$^\circ$ & 15.5--18.2$^\circ$ & 5.2\,cm \\
Circle, unannounced & 30/90/150/270 & 0.34\,s & 8.1$^\circ$ & 6.0--7.3$^\circ$ & 6.5\,cm \\
\bottomrule
\end{tabularx}
\end{table}

Every detach landed with the three remaining UAVs still carrying the ring. The freed UAV stepped clear,
but its own landing did not finish: its touchdown test fired 11--14\,cm above its resting height, its
reference froze there, and it rose to 0.34--0.38\,m and hovered until the fleet's disarm at the end of
the flight dropped it about 0.3\,m onto the mat, in all four flights. Its landing has since been changed to
push it down to its own resting height and then disarm it alone (in the twin; on the rig next). The
unannounced detaches were found 0.33--0.42\,s after the magnet was switched
off, which includes the magnet's own release time. The announced detach tilted the ring about twice
as far as in the twin, while the unannounced detaches stayed within the twin's range. In the circle the
remaining three carried on along the path without a pause.

\paragraph{A constant offset.} All four flights held the ring 5--6.5\,cm to one side of its
reference, against 1.9\,cm in the carry flights of the day before. The detach flights ran the
planner's reference height integral, because the model integral was not yet allowed on a fleet that
changes size; the carry flights ran the model integral. In the circle the offset made the path error
6.8\,cm; with the circle's mean offset and the lag removed, the remaining error was 2.4\,cm. The
model integral is now kept through a resize: in the twin it reduced the hover offset of the same
detach configurations from about 3\,cm to 2\,cm, and it is flown on the rig next.
